\documentclass[12pt,letterpaper]{article}

\usepackage[utf8]{inputenc}
\usepackage[T1]{fontenc}
\usepackage[spanish,es-tabla]{babel}
\usepackage{geometry}
\usepackage{hyperref}
\usepackage{longtable}
\usepackage{tabularx}
\usepackage{array}
\usepackage{enumitem}
\usepackage{xcolor}
\usepackage{listings}
\usepackage{booktabs}
\usepackage{microtype}
\usepackage{float}
\usepackage{graphicx}
\usepackage{amsmath}

\geometry{margin=2.5cm}
\setlength{\parskip}{0.55em}
\setlength{\parindent}{0pt}
\emergencystretch=2em

\hypersetup{
    colorlinks=true,
    linkcolor=blue,
    urlcolor=blue,
    citecolor=blue
}

\definecolor{codegray}{gray}{0.96}
\definecolor{keywordblue}{rgb}{0.05,0.19,0.45}

\lstdefinelanguage{HULK}{
    keywords={function,define,def,type,inherits,new,self,base,let,in,true,false,if,elif,else,while,for,is,as,protocol,interface,extends,match,case,default,print},
    sensitive=true,
    comment=[l]{//},
    string=[b]"
}

\lstdefinelanguage{BYTE}{
    keywords={PUSH_CONST,LOAD,STORE,ASSIGN,POP,BEGIN_SCOPE,POP_SCOPE,ADD,SUB,MUL,DIV,POW,MOD,NEG,NOT,AND,OR,CONCAT,CMP_EQ,CMP_NEQ,CMP_LT,CMP_GT,CMP_LE,CMP_GE,JUMP,JUMP_IF_FALSE,JUMP_IF_TRUE,LABEL,RETURN,HALT,CALL,NEW,GET_ATTR,SET_ATTR,SELF,METHOD_CALL,BASE_CALL,IS,AS,NEW_VECTOR,VECTOR_PUSH,VECTOR_INDEX,VECTOR_STORE,SIZE,ITER_NEXT,ITER_CURRENT,RANGE},
    sensitive=true,
    comment=[l]{\#},
    string=[b]"
}

\lstset{
    basicstyle=\ttfamily\small,
    keywordstyle=\bfseries\color{keywordblue},
    backgroundcolor=\color{codegray},
    breaklines=true,
    frame=single,
    columns=fullflexible,
    keepspaces=true,
    showstringspaces=false,
    literate={á}{{\'a}}1 {é}{{\'e}}1 {í}{{\'i}}1 {ó}{{\'o}}1 {ú}{{\'u}}1 {Á}{{\'A}}1 {É}{{\'E}}1 {Í}{{\'I}}1 {Ó}{{\'O}}1 {Ú}{{\'U}}1 {ñ}{{\~n}}1 {Ñ}{{\~N}}1
}

\newcommand{\repo}[1]{\texttt{\detokenize{#1}}}
\newcommand{\hulk}{\textsc{HULK}}

\title{Informe Final del Compilador de \hulk\\
\large Frontend con Flex/Bison, análisis semántico, generación de bytecode,
máquina virtual de pila y extensión de lenguaje mediante macros}

\author{
Ronald Alfonso Pérez , Abraham Rey Sánchez Amador, Rodrigo Meseros González\\[0.5em]
\textit{Proyecto de Compilación / Lenguajes de Programación}\\
\textit{Universidad de La Habana}
}

\date{22 de junio de 2026}

\begin{document}

\maketitle

\begin{abstract}
Este informe presenta el diseño y la implementación de un compilador para \hulk{}
construido en C++ como un sistema por fases: análisis léxico generado con Flex,
análisis sintáctico generado con Bison en modo GLR, construcción de un AST,
una fase de expansión de macros y desazucarado, análisis semántico en varias
pasadas, generación de un bytecode propio y una máquina virtual de pila que
ejecuta dicho bytecode. El énfasis del documento no está solo en enumerar
componentes del repositorio, sino en justificar las decisiones de diseño tomadas
en cada frontera del compilador: por qué se usan generadores (Flex/Bison) en
lugar de un analizador manual, por qué el parser se declara GLR, por qué las
macros y \repo{match} se resuelven como transformaciones AST$\to$AST antes de la
semántica, por qué se define un bytecode propio en lugar de emitir ensamblador o
LLVM, y por qué la ejecución se delega a una máquina virtual de pila enlazada a
través de un artefacto intermedio. Desde la perspectiva de Lenguajes de
Programación, la extensión propia principal es un \emph{sistema de macros
sintácticas} con parámetros por valor y de bloque, higiene mediante
\repo{$}-placeholders y expansión recursiva acotada, complementado por la
construcción \repo{match} con desazucarado a primitivas existentes. El objetivo
del informe es presentar de forma formal el alcance técnico del compilador, las
garantías estáticas que ofrece, los compromisos asumidos y las limitaciones que
quedan como trabajo futuro. El sistema supera el 100\,\% de la batería de pruebas
de entrega final (115/115).
\end{abstract}

\tableofcontents
\newpage

% ============================================================
\section{Introducción}

Construir un compilador para \hulk{} no consiste solamente en transformar texto
de entrada en otro texto de salida. El objetivo central es preservar el
significado de un programa a través de varias capas de representación. Un
programa fuente con expresiones, funciones, objetos, herencia, protocolos,
vectores, iterables, lambdas y una extensión como las macros debe convertirse
gradualmente en estructuras internas verificables y, finalmente, en una forma
ejecutable.

El sistema implementa una arquitectura por fases. Primero, el frontend reconoce
tokens y estructura sintáctica; después, el AST elimina detalles de la gramática
concreta; luego, una fase de expansión transforma macros y \repo{match} en
construcciones del núcleo del lenguaje; a continuación, la fase semántica
resuelve nombres, ámbitos, tipos, atributos, métodos, \repo{self}, \repo{base} y
builtins, y transpila las lambdas a tipos funtores; más adelante, el generador de
código produce un bytecode propio; finalmente, ese bytecode se serializa y se
ejecuta sobre una máquina virtual de pila.

El compilador se materializa en dos ejecutables construidos con CMake:
\repo{hulk}, que realiza todas las fases de compilación hasta producir el
bytecode, y \repo{hulk-vm}, que ejecuta el bytecode serializado. Esta separación
es importante desde el punto de vista académico: el sistema no se comporta como
una caja negra, sino como un conjunto de transformaciones observables, donde el
artefacto intermedio entre compilador y máquina virtual puede inspeccionarse.

Este informe no se organiza como una descripción cronológica del repositorio ni
como una lista de archivos. Su estructura sigue una idea de diseño: en cada fase
se explica qué problema debía resolverse, qué alternativas existían, qué decisión
se adoptó y qué consecuencias tiene esa decisión para las fases posteriores. La
arquitectura en C++ se presenta como soporte físico de la implementación; el
centro del análisis son las decisiones de compilación y de diseño de lenguaje.

Desde Lenguajes de Programación, el aporte propio principal es el sistema de
macros. Esta construcción introduce una forma de metaprogramación a nivel de
sintaxis: el programador define plantillas de código reutilizable que el
compilador expande antes de analizar el programa. Su relevancia no es únicamente
sintáctica: modifica el AST, introduce parámetros por valor y de bloque, requiere
una estrategia de higiene para evitar captura de variables y exige control de
recursión durante la expansión. La construcción \repo{match} acompaña a las
macros como azúcar sintáctica para clasificación por tipo dinámico, resuelta
mediante desazucarado a \repo{let}, \repo{if} e \repo{is}.

% ============================================================
\section{Alcance del compilador}

\subsection{Propósito general}

El propósito del compilador es transformar programas escritos en \hulk{} en un
bytecode propio y, posteriormente, ejecutarlo sobre una máquina virtual de pila.
Para lograrlo, el sistema no realiza una traducción directa desde el código
fuente hasta instrucciones de máquina, sino una sucesión de representaciones con
responsabilidades separadas: tokens, AST, AST expandido (sin macros ni
\repo{match}), AST verificado y anotado, bytecode y, finalmente, ejecución.

Esta separación permite que cada fase tenga un contrato claro. El frontend debe
reconocer la estructura del programa; la expansión debe eliminar las
construcciones de azúcar; la semántica debe rechazar programas inválidos y anotar
tipos; el generador debe producir bytecode correcto; y la máquina virtual debe
ejecutar ese bytecode con un modelo de valores y objetos bien definido.

\subsection{Lenguaje soportado}

El compilador cubre el núcleo expresivo de \hulk{} y varios mecanismos avanzados.
El lenguaje soportado incluye literales numéricos, booleanos y de cadena;
operadores aritméticos, booleanos, de comparación y de concatenación (\repo{@},
\repo{@@}); bloques; variables con \repo{let}; asignación destructiva
(\repo{:=}); condicionales con \repo{if}/\repo{elif}/\repo{else}; ciclos
\repo{while} y \repo{for}; funciones globales; llamadas recursivas; objetos con
atributos y métodos; herencia simple; despacho dinámico de métodos; \repo{self};
\repo{base}; pruebas dinámicas de tipo con \repo{is}; \emph{casts} con \repo{as};
protocolos estructurales; vectores con literales, indexación y comprensiones;
iterables y \repo{range}; lambdas y \emph{closures}; macros con \repo{define}; y
expresiones \repo{match}.

La expresión final del archivo funciona como punto de entrada del programa. Este
rasgo obliga a tratar condicionales, ciclos, bloques, llamadas y \repo{match}
como expresiones con tipo y valor resultante, no solo como instrucciones. Por
tanto, el análisis semántico calcula tipos para expresiones compuestas y la
generación de código preserva el valor producido por cada construcción.

\begin{lstlisting}[language=HULK,caption={Programa con funciones, objetos, herencia y despacho dinámico}]
function square(x: Number): Number => x * x;

type Shape() {
    area(): Number => 0;
}

type Circle(r: Number) inherits Shape {
    r: Number = r;
    area(): Number => 3 * square(self.r);
}

let s: Shape = new Circle(4) in print(s.area());
\end{lstlisting}

\subsection{Alcance técnico}

El compilador principal está implementado en C++ y se construye con CMake usando
Flex para el analizador léxico y Bison para el analizador sintáctico. La
estructura física del repositorio refleja las fases del compilador. Sin embargo,
esta división no es el objetivo central del diseño, sino un medio para mantener
límites claros entre responsabilidades.

\begin{center}
\small
\begin{tabularx}{\textwidth}{p{4.2cm}X}
\toprule
\textbf{Componente} & \textbf{Responsabilidad principal} \\
\midrule
\repo{src/lexer.l} & Especificación Flex del analizador léxico; tokens, literales, comentarios, seguimiento de línea/columna y errores léxicos. \\
\repo{src/parser.y} & Gramática Bison (GLR); precedencias, estructura de declaraciones y expresiones, construcción directa del AST. \\
\repo{include/ast.hpp} & Jerarquía de nodos del AST y patrón \emph{visitor}. \\
\repo{src/macro_expander.cpp} & Expansión de macros y desazucarado de \repo{match} (transformación AST$\to$AST). \\
\repo{src/semantic_analyzer.cpp} & Resolución de nombres, registro de tipos, chequeo semántico, inferencia acotada y transpilación de lambdas. \\
\repo{src/type_table.cpp} & Tabla de tipos, subtipado, conformidad de protocolos y tipos sintéticos (\repo{Vector}, \repo{Iterable}, funtores). \\
\repo{src/code_generator.cpp} & Generación de bytecode desde el AST verificado. \\
\repo{include/bytecode.hpp} & Definición del conjunto de instrucciones y del programa de bytecode. \\
\repo{src/serialize.cpp} & Serialización/deserialización binaria del bytecode. \\
\repo{src/vm.cpp} & Máquina virtual de pila que ejecuta el bytecode. \\
\repo{src/main.cpp} / \repo{src/vm_main.cpp} & Controladores de \repo{hulk} y \repo{hulk-vm}. \\
\bottomrule
\end{tabularx}
\end{center}

\subsection{Pipeline de alto nivel}

El pipeline completo puede resumirse de la siguiente manera:

\begin{lstlisting}[language=BYTE,caption={Pipeline conceptual del compilador}]
archivo .hulk
  -> lexer Flex                      # tokens con linea/columna
  -> parser Bison GLR                # AST (Program*)
  -> MacroExpander                   # expande macros, desazucara match
  -> SemanticAnalyzer                # 3 pasadas; tipa y transpila lambdas
  -> CodeGenerator                   # BytecodeProgram (codigo + constantes + tablas)
  -> serialize -> compiled.asm       # bytecode binario (magic "HULK")
  -> ./output (script bash)          # ejecuta ./hulk-vm compiled.asm
  -> hulk-vm                         # deserializa y ejecuta en la VM de pila
\end{lstlisting}

El diseño se apoya en una idea clásica de compilación: no intentar resolver todos
los problemas en una sola representación. El AST es cómodo para semántica, pero
demasiado alto nivel para ejecutar; el bytecode es adecuado para una máquina
virtual, pero demasiado bajo nivel para razonar sobre tipos. Por eso se introducen
fronteras intermedias, y las construcciones de azúcar (macros, \repo{match},
lambdas) se eliminan antes de llegar al bytecode.

El controlador \repo{hulk} distingue las fases mediante \emph{códigos de salida},
lo que permite su uso en correctores automáticos:

\begin{center}
\small
\begin{tabularx}{\textwidth}{p{3.0cm}p{2.2cm}X}
\toprule
\textbf{Fase} & \textbf{Código} & \textbf{Significado} \\
\midrule
Léxico & 1 & Error léxico (\repo{lexical_errors_count > 0}). \\
Sintáctico / Macro & 2 & Error sintáctico o de expansión de macros. \\
Semántico & 3 & Error de análisis semántico. \\
Éxito & 0 & Se generó \repo{compiled.asm} y \repo{./output}. \\
\bottomrule
\end{tabularx}
\end{center}

% ============================================================
\section{Frontend}

El frontend transforma una secuencia de caracteres en una representación
estructurada del programa. Esta fase resuelve dos problemas diferentes. El
primero es léxico: identificar palabras clave, operadores, delimitadores,
identificadores, literales, comentarios y espacios. El segundo es sintáctico:
reconocer cómo esos tokens forman declaraciones, expresiones, tipos, protocolos,
bloques y patrones.

\subsection{Decisión léxica: lexer generado con Flex}

El analizador léxico se especifica con Flex en \repo{src/lexer.l}. Esta decisión
es adecuada porque el conjunto léxico de \hulk{} no se reduce a identificadores y
números. El lenguaje contiene palabras clave (\repo{function}, \repo{define},
\repo{type}, \repo{inherits}, \repo{new}, \repo{self}, \repo{base}, \repo{let},
\repo{if}, \repo{elif}, \repo{else}, \repo{while}, \repo{for}, \repo{protocol},
\repo{interface}, \repo{extends}, \repo{is}, \repo{as}, \repo{match},
\repo{case}, \repo{default}); operadores de varios caracteres (\repo{@@},
\repo{@}, \repo{==}, \repo{!=}, \repo{<=}, \repo{>=}, \repo{:=}, \repo{=>},
\repo{->}); delimitadores; literales numéricos (enteros y flotantes, incluyendo
la forma \repo{.5}); literales de cadena con secuencias de escape; comentarios de
línea \repo{//}; e identificadores, incluyendo identificadores con prefijo
\repo{$} usados por el sistema de macros.

Un lexer manual daría control fino sobre cada caso, pero haría que la definición
del lenguaje quedara dispersa en código imperativo. Al usar Flex, la
especificación léxica se vuelve declarativa: las reglas que describen tokens se
mantienen separadas del código que los consume. Desde el punto de vista teórico,
esta decisión conecta la implementación con la idea de autómatas finitos y
expresiones regulares para el reconocimiento léxico.

El lexer mantiene un seguimiento explícito de posición mediante las variables
globales \repo{line_number} y \repo{column_number}: cada token incrementa la
columna según su longitud y los saltos de línea reinician la columna. Esta
información se propaga al parser y al AST, lo que permite que todos los
diagnósticos posteriores indiquen \repo{(línea,columna)}.

Los errores léxicos se reportan con un formato uniforme y se contabilizan en
\repo{lexical_errors_count}; entre ellos están la secuencia de escape inválida, la
cadena sin cerrar (en fin de línea o de archivo) y el carácter inesperado:

\begin{lstlisting}[language=BYTE,caption={Formato de error léxico}]
(linea,columna) LEXICAL: unexpected character 'X'
(linea,columna) LEXICAL: unterminated string literal
\end{lstlisting}

\subsection{Decisión sintáctica: Bison en modo GLR}

El analizador sintáctico se especifica con Bison en \repo{src/parser.y} y se
declara \repo{\%glr-parser}. La elección de GLR (Generalized LR) en lugar de LALR
puro es deliberada: \hulk{} es un lenguaje basado en expresiones con construcciones
que generan conflictos de desplazamiento/reducción difíciles de eliminar por
completo en LALR (por ejemplo, lambdas como operandos, o construcciones de control
que pueden aparecer como operando derecho de un binario). GLR permite que el
generador conserve varias derivaciones cuando una decisión no puede resolverse
localmente, y resuelve la mayoría de los conflictos mediante las declaraciones de
precedencia y asociatividad.

La gramática define una escalera de precedencia explícita, de menor a mayor poder
de enlace: \repo{in} y \repo{else}; asignación \repo{:=}; flechas \repo{=>} y
\repo{->}; disyunción \repo{|}; conjunción \repo{&}; negación \repo{!};
comparaciones y \repo{is}/\repo{as} (no asociativas, lo que impide encadenarlas);
concatenación \repo{@}/\repo{@@}; suma y resta; multiplicación, división y módulo;
potencia \repo{^}; y negación unaria. Esta declaración hace explícito el diseño de
precedencias y evita multiplicar reglas intermedias.

La estructura general distingue declaraciones de nivel superior (funciones, tipos,
protocolos y macros) de la expresión global final. Las expresiones se organizan en
una jerarquía descendente (\repo{assign\_expr}, \repo{or\_expr}, \dots,
\repo{primary\_expr}) coherente con la escalera de precedencia.

\paragraph{Construcciones de control como operando derecho.} Una decisión de
diseño reciente y relevante es permitir que \repo{if}, \repo{while}, \repo{for} y
\repo{match} aparezcan como operando \emph{derecho} de un operador binario, de modo
que expresiones como
\begin{lstlisting}[language=HULK]
evens := evens + if (i % 2 == 0) 1 else 0;
\end{lstlisting}
sean válidas. La solución adoptada agrupa estas construcciones en un no terminal
\repo{flow\_expr} que \emph{no} es un \repo{primary\_expr} (para que no puedan ser
operando \emph{izquierdo}, lo que volvería ambiguo \repo{if (c) x else a + b}) y se
añade explícitamente como operando derecho de cada nivel binario. Como cada
construcción comienza con una palabra clave distinta, el reconocimiento es
determinista y el cuerpo de la construcción se extiende de forma \emph{greedy}
hacia la derecha, preservando la semántica esperada. Esta decisión fue necesaria
para superar el último caso de la batería de pruebas.

Los errores sintácticos se reportan con un formato uniforme que incluye posición y
el token problemático:

\begin{lstlisting}[language=BYTE,caption={Formato de error sintáctico}]
(linea,columna) SYNTACTIC: <mensaje> near '<token>'
\end{lstlisting}

\subsection{AST como frontera entre sintaxis y semántica}

El AST (\repo{include/ast.hpp}) es la primera representación estable del lenguaje.
El parser construye directamente nodos de AST; no se mantiene un CST explícito
separado, sino estructuras temporales de \emph{parsing} que se convierten en
vectores de punteros únicos antes de armar los nodos finales. Esta separación evita
que la fase semántica dependa de detalles accidentales de la gramática.

El AST emplea un patrón \emph{visitor} genérico (\repo{Visitor<T>}) con un método
\repo{visit} por cada tipo de nodo. Existen alrededor de treinta tipos de nodos,
agrupados en: nodos de expresión (literales, \repo{BinaryExpr}, \repo{UnaryExpr},
\repo{VarRef}, \repo{AssignExpr}, \repo{BlockExpr}, \repo{IfExpr},
\repo{WhileExpr}, \repo{ForExpr}, \repo{LetExpr}, \repo{LambdaExpr},
\repo{FuncCall}, \repo{MethodCall}, \repo{MemberAccess}, \repo{NewExpr},
\repo{VectorLiteral}, comprensiones, \repo{VectorIndex}, \repo{ArrayNewExpr},
\repo{ArrayAssignExpr}, \repo{IsExpr}, \repo{AsExpr}, \repo{SelfRef},
\repo{BaseCall}); nodos de declaración (\repo{FuncDef}, \repo{TypeDef},
\repo{AttributeDef}, \repo{MethodDef}, \repo{ProtocolDef}, \repo{MacroDef}); y el
nodo raíz \repo{Program}, que agrupa tipos, protocolos, funciones, macros y la
expresión global.

Dos nodos tienen un comportamiento especial que documenta una decisión de diseño:
\repo{MacroInvoke} y \repo{MatchExpr} \emph{lanzan una excepción} si se les invoca
\repo{accept()}. Esto garantiza, por construcción, que ambos deben haber sido
eliminados por la fase de expansión antes de llegar a la semántica o a la
generación de código. El compilador no permite, por diseño, que una macro o un
\repo{match} sin expandir alcancen fases posteriores.

% ============================================================
\section{Análisis semántico}

Después del frontend y de la expansión, el compilador conoce la estructura del
programa, pero todavía no sabe si esa estructura tiene sentido. La fase semántica
(\repo{src/semantic_analyzer.cpp}) responde preguntas que la gramática no puede
resolver: si una variable está definida, si una llamada tiene la aridad correcta,
si un tipo existe, si una condición es booleana, si un método pertenece al
receptor, si un \emph{cast} es admisible, si una redefinición respeta la firma del
padre o si un tipo satisface un protocolo.

\subsection{Análisis en varias pasadas}

El análisis se organiza en tres pasadas secuenciales sobre el AST. Esta decisión
responde a una necesidad del lenguaje: las declaraciones globales pueden referirse
entre sí, las funciones pueden ser recursivas y la herencia requiere conocer la
jerarquía completa antes de validar cuerpos.

\begin{enumerate}[leftmargin=1.6em]
\item \textbf{Registro de tipos.} Recolecta todas las definiciones de tipo,
valida ausencia de tipos, atributos, métodos y parámetros duplicados, registra
protocolos, comprueba que no existan ciclos de herencia y registra los tipos en
orden topológico (los padres antes que los hijos).
\item \textbf{Registro de funciones.} Recolecta las funciones globales, valida que
no haya nombres ni parámetros duplicados y almacena sus firmas.
\item \textbf{Chequeo de tipos.} Verifica cuerpos de métodos, inicializadores de
atributos, cuerpos de funciones y la expresión global; acumula errores y, si los
hay, los imprime y aborta con código de salida 3.
\end{enumerate}

\subsection{Tabla de tipos, subtipado y protocolos}

La tabla de tipos (\repo{src/type_table.cpp}) registra los tipos primitivos
\repo{Number}, \repo{String}, \repo{Boolean} y el tipo raíz \repo{Object}, además
de tipos sintéticos generados bajo demanda: \repo{Vector<T>}, \repo{Iterable<T>} y
tipos funtor (\repo{\_FunctorType\_N}) para valores invocables. Los tipos
primitivos son \emph{finales}: heredar de ellos es un error.

El subtipado nominal recorre la cadena de padres: todo tipo conforma con
\repo{Object}, y un tipo conforma con un ancestro si aparece en su cadena de
herencia. Los contenedores son covariantes (\repo{Vector<T> <= Vector<U>} si
\repo{T <= U}) y un \repo{Vector<T>} conforma con \repo{Iterable<T>}.

Para protocolos, el sistema valida conformidad \emph{estructural}: un tipo puede
usarse donde se espera un protocolo si posee todos los métodos requeridos con
firmas compatibles, comprobando además reglas de varianza (contravarianza en
parámetros y covarianza en el tipo de retorno). Esta decisión introduce una forma
de polimorfismo basada en comportamiento, no en nombres.

\subsection{Inferencia acotada}

\hulk{} permite omitir anotaciones de tipo en varios contextos, por lo que el
compilador realiza una inferencia \emph{acotada} (no un sistema general de
inferencia). Las decisiones principales son: los parámetros sin anotar adoptan por
defecto el tipo \repo{Number} (heurística justificada porque \hulk{} no tiene
sobrecarga de operadores y un parámetro usado aritméticamente solo puede ser
numérico); el tipo de un atributo sin anotar se infiere de su inicializador; el
tipo de una literal de vector es el ancestro común más bajo (LCA) de sus
elementos; el tipo de un \repo{if} es el LCA de sus ramas; y el retorno de una
función puede inferirse de su cuerpo. Se trata de una inferencia local de una sola
pasada, no de un punto fijo iterativo: un compromiso entre expresividad y
simplicidad.

\subsection{Resolución de nombres, \repo{self} y \repo{base}}

La resolución de nombres usa una tabla de símbolos basada en una pila de ámbitos:
ámbito global (con las constantes \repo{PI} y \repo{E} predefinidas), ámbitos de
función, ámbitos de \repo{let} y ámbitos de método. Las búsquedas van del ámbito
más interno al más externo. \repo{self} solo es válido dentro de métodos y su tipo
es el del tipo contenedor; no es un objetivo de asignación válido. \repo{base()}
solo es válido dentro de métodos y se resuelve contra el método homónimo del tipo
padre, lo que distingue la llamada al ancestro del despacho dinámico ordinario.

\subsection{Transpilación de lambdas a funtores}

Una decisión central de esta fase es que las lambdas no se tratan como un valor
primitivo del runtime, sino que se \emph{transpilan} a un tipo funtor. Al visitar
una \repo{LambdaExpr}, el analizador detecta las variables libres capturadas del
entorno, genera un \repo{TypeDef} sintético (\repo{\_Lambda\_N}) con un atributo
por cada captura y un método \repo{invoke} cuyo cuerpo es el de la lambda, y anota
en el nodo el nombre del tipo generado y la lista de capturas. El cuerpo original
se envuelve en un \repo{let} que reenlaza cada variable capturada a
\repo{self.\_var}, de modo que el código de la lambda puede seguir refiriéndose a
los nombres originales. Esta estrategia reduce la complejidad del backend: una
lambda se convierte en \repo{new \_Lambda\_N(capturas)} y su invocación en una
llamada al método \repo{invoke}.

\subsection{Errores semánticos diferenciados}

Una fortaleza de la fase semántica es que distingue errores de naturaleza
diferente, todos con el formato \repo{(línea,columna) SEMANTIC: mensaje}. Entre los
casos cubiertos aparecen: tipos, funciones, métodos, atributos y parámetros
duplicados; símbolos y métodos indefinidos; objetivos de asignación inválidos;
aridades incorrectas en funciones, métodos, constructores y \repo{base()};
herencia circular; herencia desde primitivos; padre inexistente; \emph{override}
con aridad, firma o retorno incompatibles; atributos privados accedidos desde
fuera de su tipo; condiciones no booleanas en \repo{if}, \repo{while} y filtros de
comprensión; \repo{for} sobre algo no iterable; operandos inválidos en operadores
aritméticos, booleanos, de comparación y de concatenación; \emph{casts}
inadmisibles; e inferencia no resuelta. En una defensa académica esto es esencial:
un compilador no se evalúa únicamente por los programas que ejecuta, sino también
por los programas inválidos que detecta antes de generar código.

\subsection{Funciones y constantes incorporadas}

El sistema predefine las funciones \repo{print}, \repo{sqrt}, \repo{sin},
\repo{cos}, \repo{exp}, \repo{log} (con base), \repo{rand} y \repo{range}, así como
las constantes \repo{PI} y \repo{E}. La función \repo{range} produce un objeto que
conforma con el protocolo \repo{Iterable}, lo que permite recorrerlo con
\repo{for}.

% ============================================================
\section{Generación de bytecode}

\subsection{Motivación de un bytecode propio}

Una decisión central del sistema es no emitir ensamblador nativo ni LLVM, sino un
\emph{bytecode propio} ejecutado por una máquina virtual. Esta elección responde al
alcance del proyecto: generar ensamblador exigiría manejar convenciones de llamada,
registros, pila, alineación y detalles de plataforma; integrar LLVM añadiría una
dependencia y un nivel de abstracción que no es necesario para los objetivos del
curso. Un bytecode propio, en cambio, permite concentrar el trabajo en la
traducción semántica de \hulk{} y mantener una representación legible y depurable,
cercana a una máquina de pila.

\subsection{Modelo de pila y conjunto de instrucciones}

El bytecode (\repo{include/bytecode.hpp}) está orientado a una \emph{máquina de
pila}: los operandos se empujan y se consumen de una pila de valores. El conjunto
de instrucciones cubre movimiento de datos y ámbitos, aritmética, lógica,
concatenación, comparación, control de flujo, llamadas, modelo de objetos,
vectores, iteración y rangos.

\begin{center}
\small
\begin{tabularx}{\textwidth}{p{3.6cm}X}
\toprule
\textbf{Grupo} & \textbf{Instrucciones representativas} \\
\midrule
Datos y ámbitos & \repo{PUSH\_CONST}, \repo{LOAD}, \repo{STORE}, \repo{ASSIGN}, \repo{POP}, \repo{BEGIN\_SCOPE}, \repo{POP\_SCOPE}. \\
Aritmética y lógica & \repo{ADD}, \repo{SUB}, \repo{MUL}, \repo{DIV}, \repo{POW}, \repo{MOD}, \repo{NEG}, \repo{NOT}, \repo{AND}, \repo{OR}. \\
Comparación y cadenas & \repo{CMP\_EQ}, \repo{CMP\_NEQ}, \repo{CMP\_LT}, \repo{CMP\_GT}, \repo{CMP\_LE}, \repo{CMP\_GE}, \repo{CONCAT}. \\
Control de flujo & \repo{JUMP}, \repo{JUMP\_IF\_FALSE}, \repo{JUMP\_IF\_TRUE}, \repo{LABEL}, \repo{RETURN}, \repo{HALT}. \\
Funciones & \repo{CALL}. \\
Objetos & \repo{NEW}, \repo{GET\_ATTR}, \repo{SET\_ATTR}, \repo{SELF}, \repo{METHOD\_CALL}, \repo{BASE\_CALL}. \\
Tipos dinámicos & \repo{IS}, \repo{AS}. \\
Vectores e iteración & \repo{NEW\_VECTOR}, \repo{VECTOR\_PUSH}, \repo{VECTOR\_INDEX}, \repo{VECTOR\_STORE}, \repo{SIZE}, \repo{ITER\_NEXT}, \repo{ITER\_CURRENT}, \repo{RANGE}. \\
\bottomrule
\end{tabularx}
\end{center}

Un \repo{BytecodeProgram} contiene cuatro secciones: la lista de instrucciones, una
tabla de constantes (con deduplicación), una tabla de funciones (que mapea nombres
a índices de instrucción, incluyendo métodos como \repo{Tipo.metodo} y
constructores como \repo{Tipo.\_\_init\_\_}) y una tabla de ancestros por tipo,
usada para el despacho dinámico en tiempo de ejecución.

\subsection{Lowering de objetos, métodos y \repo{base}}

El generador de código (\repo{src/code_generator.cpp}) construye, para cada
\repo{TypeDef}, un método constructor \repo{Tipo.\_\_init\_\_}. Si el tipo tiene
padre, el constructor comienza con un \repo{BASE\_CALL} al \repo{\_\_init\_\_} del
ancestro; luego inicializa los atributos a partir de los parámetros del constructor
y de los inicializadores, y devuelve \repo{self}. La instanciación \repo{new T(...)}
emite \repo{NEW T}, evalúa los argumentos y llama \repo{\_\_init\_\_}.

El despacho dinámico se resuelve en tiempo de ejecución: \repo{METHOD\_CALL}
construye el símbolo \repo{tipo.metodo} y, si no lo encuentra, recorre la cadena de
ancestros del objeto hasta hallar la primera implementación. Esto logra el efecto
de una tabla virtual sin almacenar punteros de vtable, lo que simplifica la
serialización. La llamada \repo{base(...)} se compila como \repo{BASE\_CALL}
estático contra el método del ancestro, evitando que vuelva a invocarse el
\emph{override} del hijo.

\subsection{Lowering de \repo{self}, vectores, ciclos y lambdas}

El \repo{self} implícito de los métodos se trata como un valor disponible mediante
la instrucción \repo{SELF}. Los vectores literales se construyen con
\repo{NEW\_VECTOR}; la indexación usa \repo{VECTOR\_INDEX}; la asignación por índice,
\repo{VECTOR\_STORE}; y la longitud, \repo{SIZE}. Los ciclos \repo{for} y las
comprensiones se bajan a un patrón basado en \repo{ITER\_NEXT}/\repo{ITER\_CURRENT}
sobre el iterable. Las lambdas, ya transpiladas a tipos funtor por la semántica, se
emiten como \repo{NEW \_Lambda\_N}, seguido de la carga de las variables capturadas
y la llamada a \repo{\_\_init\_\_}; su invocación se compila como una llamada al
método \repo{invoke}.

\begin{lstlisting}[language=BYTE,caption={Esquema de bytecode para un \texttt{for}}]
STORE  iter_var          # guarda el iterable
LABEL  loop_start
LOAD   iter_var
ITER_NEXT                # empuja Boolean: hay siguiente?
JUMP_IF_FALSE loop_end
LOAD   iter_var
ITER_CURRENT             # empuja el elemento actual
STORE  x                 # variable del for
...  cuerpo ...
JUMP   loop_start
LABEL  loop_end
\end{lstlisting}

\subsection{Serialización}

El programa de bytecode se serializa en formato \emph{binario}
(\repo{src/serialize.cpp}) a un archivo \repo{compiled.asm}. El formato comienza con
un número mágico (\repo{0x48554C4B}, las letras ``HULK'') y una versión, seguidos de
las secciones de código, constantes, tabla de funciones y tabla de ancestros. La
máquina virtual deserializa este archivo para reconstruir el \repo{BytecodeProgram}.

% ============================================================
\section{Máquina virtual de ejecución}

\subsection{Arquitectura}

La máquina virtual (\repo{src/vm.cpp}) es un intérprete de pila. Su estado consiste
en una pila de operandos, un puntero de instrucción, una pila de marcos de llamada
(cada uno con dirección de retorno, entorno y \repo{self}) y una cadena de entornos
para el ámbito de variables. El bucle principal recorre las instrucciones y las
ejecuta mediante un \emph{switch} sobre el código de operación.

\subsection{Modelo de valores y objetos}

El valor de runtime (\repo{include/value.hpp}) es una unión etiquetada con los
casos \repo{Number} (un \repo{double}), \repo{String}, \repo{Boolean},
\repo{Object}, \repo{Vector} y \repo{Null}. Los objetos y vectores se gestionan con
punteros compartidos. Un \repo{HulkObject} (\repo{include/hulk_object.hpp}) almacena
su nombre de tipo, su cadena de ancestros y un mapa de atributos; el despacho de
métodos usa esa cadena de ancestros. Los iterables del runtime, como
\repo{HulkVector} y \repo{HulkRange}, exponen \repo{next()} y \repo{current()}, lo
que permite que \repo{ITER\_NEXT} e \repo{ITER\_CURRENT} operen uniformemente sobre
vectores, rangos y objetos definidos por el usuario que conformen con
\repo{Iterable}.

\subsection{Llamadas, despacho y builtins}

Una instrucción \repo{CALL} busca primero en la tabla de funciones del programa; si
no la encuentra, busca en el entorno global una función nativa (los builtins
\repo{print}, \repo{sqrt}, \repo{range}, etc., implementados como objetos de función
nativa). \repo{METHOD\_CALL} implementa el despacho dinámico recorriendo la cadena
de ancestros del receptor, mientras que \repo{BASE\_CALL} salta directamente a la
implementación del ancestro. La máquina virtual lanza errores de ejecución ante
situaciones como división por cero, índices fuera de rango, \emph{casts} inválidos o
métodos inexistentes.

\subsection{Por qué la VM no es ``otro backend''}

Conviene precisar que la máquina virtual no es un backend alternativo al generador
de bytecode, sino su consumidor natural. El compilador \repo{hulk} produce el
bytecode y un pequeño \emph{script} \repo{./output} que invoca \repo{hulk-vm} sobre
\repo{compiled.asm}. Así, el flujo de artefactos es claro: una sola tubería de
compilación seguida de una etapa de ejecución desacoplada.

% ============================================================
\section{Extensión propia de lenguaje: macros y \repo{match}}

\subsection{Problema de diseño}

\hulk{} ofrece funciones, condicionales y operadores, pero no un mecanismo para
\emph{abstraer patrones de código} que el programador repite con pequeñas
variaciones, ni para generar código a partir de plantillas. La extensión propia
principal de este proyecto es un \emph{sistema de macros sintácticas}: el
programador declara con \repo{define} (alias \repo{def}) plantillas que el
compilador expande antes del análisis semántico, sustituyendo parámetros por los
argumentos de cada invocación.

\subsection{Sintaxis y representación}

Una macro se declara como una función, pero con \repo{define}, y vive en un espacio
de nombres separado (\repo{Program.macros}). Sus parámetros pueden ser
\emph{por valor} (\repo{name}) o \emph{sintácticos} (\repo{*name}), estos últimos
destinados a recibir un bloque de código. En el AST, una declaración es un
\repo{MacroDef} (con \repo{MacroParam} que marca \repo{is\_syntactic}) y cada
invocación es un \repo{MacroInvoke}, que puede llevar un argumento de bloque.

\begin{lstlisting}[language=HULK,caption={Definición y uso de macros}]
define double(x: Number): Number -> x * 2;
define triple(x: Number): Number -> x * 3;
define six_times(x: Number): Number -> double(triple(x));

{
    if (six_times(3) == 18) print("ok") else print("fail");
}
\end{lstlisting}

\subsection{Expansión, higiene y recursión}

La expansión la realiza \repo{MacroExpander} (\repo{src/macro_expander.cpp}) como
una transformación AST$\to$AST previa a la semántica. El expansor indexa las macros
declaradas, recorre el árbol y, al encontrar un \repo{MacroInvoke} (o una llamada
cuyo nombre coincide con una macro), valida la aridad, construye un mapa de
sustitución parámetro$\to$argumento, \emph{clona} el cuerpo de la macro y sustituye.

La \emph{higiene} se resuelve con identificadores con prefijo \repo{$}: dentro del
cuerpo de una macro, un nombre como \repo{$tmp} se renombra a un identificador único
generado con un contador global (un prefijo \repo{\_\_mN\_}), de modo que distintas
expansiones no colisionan entre sí ni capturan variables del sitio de invocación.
La expansión es recursiva: el cuerpo de una macro puede invocar otras macros, que se
re-expanden, con un límite de profundidad (64) como salvaguarda frente a recursión
infinita. La existencia de pruebas dedicadas (\repo{define\_hygiene},
\repo{define\_nested}, \repo{define\_recursive\_expand}, \repo{define\_loop},
\repo{define\_conditional}, \repo{define\_chain}, \repo{define\_block}) documenta el
alcance del sistema.

\subsection{La construcción \repo{match}}

La construcción \repo{match} acompaña a las macros como azúcar sintáctica para
clasificación por \emph{tipo dinámico}. Su forma es:

\begin{lstlisting}[language=HULK,caption={Forma general de match}]
match (subject) {
    case TipoA: cuerpo_a;
    case TipoB: cuerpo_b;
    default: cuerpo_por_defecto;
}
\end{lstlisting}

El expansor desazucara \repo{match} a una combinación de primitivas ya existentes:
evalúa el \emph{scrutinee} una sola vez ligándolo a una variable sintética con
\repo{let}, y genera una cadena \repo{if}/\repo{elif}/\repo{else} en la que cada
brazo prueba \repo{variable is Tipo}; el brazo \repo{default} se convierte en el
\repo{else} final. Si no hay \repo{default}, se usa un valor por defecto como
respaldo. Esta estrategia tiene varias ventajas: no agrega instrucciones nuevas al
bytecode, reutiliza las reglas semánticas de \repo{let}, \repo{if} e \repo{is}, y
mantiene la máquina virtual independiente de una construcción de patrones.

\begin{lstlisting}[language=HULK,caption={Esquema conceptual de desazucarado de match}]
let __match0__ = subject in
    if (__match0__ is TipoA) cuerpo_a
    elif (__match0__ is TipoB) cuerpo_b
    else cuerpo_por_defecto
\end{lstlisting}

La evaluación única del \emph{scrutinee} es importante: si tuviera efectos
secundarios, reevaluarlo por cada brazo cambiaría el significado del programa; por
eso el desazucarado lo guarda primero en una variable.

\subsection{Comparación con otros lenguajes}

El sistema de macros de \hulk{} se inspira en los macros de sustitución sintáctica
presentes en la tradición de Lisp y en el preprocesador de C, pero adopta una
posición intermedia. A diferencia del preprocesador de C, opera sobre el AST y no
sobre texto, por lo que respeta la estructura del lenguaje y puede aplicar higiene;
a diferencia de los macros \repo{syntax-rules} de Scheme, su modelo de higiene es
más simple (basado en \repo{$}-placeholders y renombrado por contador) y no implementa
\emph{pattern matching} sobre la forma sintáctica de los argumentos. La construcción
\repo{match}, por su parte, recuerda al \repo{match} de Rust o al \repo{when} de
Kotlin en la idea de seleccionar por tipo dinámico, pero su alcance es deliberadamente
menor: solo patrones de tipo y un brazo por defecto, sin \emph{binding} ni
\emph{narrowing} explícito, lo que basta para el modelo de objetos con herencia
abierta de \hulk{} sin asumir tipos suma cerrados.

% ============================================================
\section{Característica adicional: comprensión de listas y generadores perezosos}

Además del sistema de macros, el compilador incorpora dos características
relacionadas con el procesamiento de secuencias: la \emph{comprensión de listas}
(construcción declarativa de vectores) y los \emph{generadores perezosos}
(secuencias producidas bajo demanda). Ambas comparten el mismo protocolo de
iteración, pero difieren en un punto esencial: la comprensión \emph{materializa}
un resultado, mientras que un generador \emph{difiere} el cálculo de sus elementos.

\subsection{Comprensión de listas}

La comprensión de listas permite construir un vector describiendo cómo se genera
cada elemento a partir de los elementos de una secuencia fuente, con un filtro
opcional. La gramática (\repo{src/parser.y}) reconoce dos formas, que producen los
nodos \repo{VectorComprehension} y \repo{VectorComprehensionFilter}:

\begin{lstlisting}[language=HULK,caption={Comprensión de listas, con y sin filtro}]
let squares = [x * x | x in range(0, 5)] in print(squares);
let evens   = [x | x in range(0, 10) if x % 2 == 0] in print(evens);
\end{lstlisting}

Semánticamente, el analizador exige que la fuente sea iterable ---un
\repo{Iterable<T>}, un \repo{Vector<T>} o un tipo que conforme con el protocolo
\repo{Iterable}--- e infiere el tipo del elemento a partir de ella; en la forma con
filtro, además verifica que la condición sea \repo{Boolean}.

El \emph{lowering} de la comprensión es \emph{ansioso} (\emph{eager}): el generador
de código (\repo{src/code_generator.cpp}) inicializa un vector vacío con
\repo{VECTOR\_INIT} y recorre la fuente con el mismo protocolo de iteración que el
\repo{for}, empujando con \repo{VECTOR\_PUSH} cada elemento generado que pase el
filtro. El resultado es un \repo{HulkVector} concreto en memoria.

\begin{lstlisting}[language=BYTE,caption={Bytecode de una comprensión con filtro}]
... evalua la fuente ...
STORE _iter
VECTOR_INIT
STORE _vec
LABEL loop_start
LOAD  _iter
ITER_NEXT                # hay siguiente?
JUMP_IF_FALSE loop_end
LOAD  _iter
ITER_CURRENT
STORE x                  # variable de la comprension
...  evalua el filtro ...
JUMP_IF_FALSE loop_start  # si el filtro es falso, no se empuja
LOAD  _vec
...  evalua el generador ...
VECTOR_PUSH
STORE _vec
JUMP loop_start
LABEL loop_end
LOAD  _vec               # el vector materializado es el valor resultante
\end{lstlisting}

\subsection{Generadores perezosos}

La fuente de una comprensión ---o de un \repo{for}--- no necesita existir
completamente en memoria. El protocolo de iteración del lenguaje es
\emph{perezoso} y basado en \emph{pull}: un iterable expone los métodos
\repo{next()} (avanza y responde si quedan elementos) y \repo{current()} (devuelve
el elemento actual), y las instrucciones \repo{ITER\_NEXT} e \repo{ITER\_CURRENT}
solicitan \emph{un} elemento a la vez. Así, los valores se calculan en el momento
en que se consumen y nunca se almacena la secuencia completa de la fuente.

El lenguaje ofrece una anotación de tipo dedicada para estas secuencias perezosas:
el sufijo \repo{*} sobre un tipo. La gramática (\repo{src/parser.y}) desazucara
\repo{T*} a \repo{Iterable<T>}, de modo que un parámetro declarado \repo{gen:
Number*} acepta cualquier generador perezoso de números. El tipo \repo{Iterable}
está predefinido como protocolo con los métodos \repo{next()} y \repo{current()}
(\repo{src/type_table.cpp}), y la función incorporada \repo{range} produce
precisamente un generador perezoso de este tipo.

Un generador perezoso puede definirse como un tipo ordinario que implemente el
protocolo: su estado interno avanza en \repo{next()} y el valor se computa en
\repo{current()}. El siguiente ejemplo ---tomado de la batería de pruebas--- define
un generador de cuadrados y lo consume mediante un parámetro de tipo \repo{Number*}:

\begin{lstlisting}[language=HULK,caption={Generador perezoso y su consumo}]
type Squares(n: Number) {
    i: Number = 0;
    limit: Number = n;
    next(): Boolean { self.i := self.i + 1; self.i <= self.limit; }
    current(): Number { self.i * self.i; }
}

function sum_gen(gen: Number*): Number {
    let s: Number = 0 in {
        for (x in gen) { s := s + x; };
        s;
    };
}

print(sum_gen(new Squares(3)));   // 1 + 4 + 9 = 14
\end{lstlisting}

En tiempo de ejecución, \repo{ITER\_NEXT} sobre un objeto distingue tres casos
(\repo{src/vm.cpp}): si es un \repo{HulkRange} llama a su \repo{next()} nativo; si es
un \repo{HulkVector} avanza su iterador interno; y si es un objeto de usuario, busca
y ejecuta su método \repo{next()} (recorriendo la cadena de ancestros), e
igualmente \repo{ITER\_CURRENT} invoca \repo{current()}. Este despacho uniforme es lo
que permite que \repo{for} y las comprensiones operen indistintamente sobre rangos,
vectores y generadores definidos por el usuario.

\subsection{Relación entre ambas características}

La distinción de diseño es clara: un \emph{generador} representa una secuencia
perezosa que puede pasarse, recorrerse y consumirse elemento a elemento sin
materializarse; una \emph{comprensión} consume una de esas secuencias (perezosa o
no) y produce un vector concreto. Dicho de otro modo, la comprensión es un
\emph{consumidor ansioso} construido sobre el \emph{protocolo perezoso} de
iteración. Esta separación permite componer ambos mundos: una comprensión puede
alimentarse de \repo{range} o de un generador de usuario, tomando perezosamente de
la fuente solo lo necesario para construir su resultado.

% ============================================================
\section{Validación y resultados}

El compilador se valida contra la batería de pruebas de entrega final, organizada
en categorías obligatorias y de bonificación. Cada caso compila un programa con
\repo{hulk}, ejecuta \repo{./output} y compara la salida con la esperada; los casos
de error verifican además el código de salida y la categoría del diagnóstico
(\repo{LEXICAL}, \repo{SYNTACTIC}, \repo{SEMANTIC}). El sistema supera el
\textbf{100\,\%} de las pruebas (\textbf{115/115}).

\begin{center}
\small
\begin{tabularx}{\textwidth}{p{4.2cm}p{2.4cm}X}
\toprule
\textbf{Categoría} & \textbf{Resultado} & \textbf{Cobertura} \\
\midrule
\repo{ok/minimal} & 20/20 & Expresiones, \repo{let}, control de flujo, funciones. \\
\repo{ok/types} & 10/10 & Tipos, atributos, métodos, retornos. \\
\repo{ok/oop} & 10/10 & Herencia, despacho dinámico, \repo{self}, \repo{base}. \\
\repo{errors/lexical} & 6/6 & Diagnósticos léxicos y código de salida 1. \\
\repo{errors/syntactic} & 10/10 & Diagnósticos sintácticos y código de salida 2. \\
\repo{errors/semantic} & 15/15 & Diagnósticos semánticos y código de salida 3. \\
\repo{ok/macros} & 8/8 & Definición, higiene, anidamiento, expansión recursiva. \\
\repo{ok/interfaces} & 6/6 & Protocolos y conformidad estructural. \\
\repo{ok/arrays} & 8/8 & Vectores, indexación, mutación, recorrido. \\
\repo{ok/lambdas} & 6/6 & Lambdas, \emph{closures}, composición. \\
\repo{ok/generators} & 6/6 & Iterables definidos por el usuario. \\
\repo{ok/extras} & 10/10 & Casos combinados, incl. control de flujo como operando. \\
\bottomrule
\end{tabularx}
\end{center}

\section{Construcción y reproducibilidad}

El proyecto se construye con CMake y requiere Flex y Bison. El objetivo de
construcción genera ambos ejecutables y los copia a la raíz del repositorio:

\begin{lstlisting}[language=BYTE,caption={Construcción}]
cmake -B build -S .
cmake --build build -j
cp build/hulk . && cp build/hulk-vm .
\end{lstlisting}

\section{Limitaciones y trabajo futuro}

El diseño asume varios compromisos conscientes. La inferencia de tipos es local y
adopta \repo{Number} como tipo por defecto de parámetros sin anotar, lo que puede
ser conservador en programas que usan parámetros no numéricos sin anotación. La
higiene de macros se basa en \repo{$}-placeholders explícitos y no en un sistema
completo de higiene automática. La construcción \repo{match} solo admite patrones de
tipo y un brazo por defecto, sin \emph{binding} ni patrones estructurales. El
despacho dinámico recorre linealmente la cadena de ancestros en lugar de usar tablas
virtuales materializadas. Como trabajo futuro se plantean: una inferencia más
expresiva, higiene automática de macros, patrones con \emph{binding} y
\emph{narrowing} en \repo{match}, y optimizaciones de bytecode (vtables, plegado de
constantes).

\section{Conclusiones}

El compilador de \hulk{} presentado implementa, de extremo a extremo, una tubería de
compilación por fases en C++: análisis léxico con Flex, análisis sintáctico GLR con
Bison, un AST con patrón \emph{visitor}, una fase de expansión que elimina macros y
\repo{match}, un análisis semántico en tres pasadas con sistema de tipos, protocolos
estructurales e inferencia acotada, la transpilación de lambdas a tipos funtor, la
generación de un bytecode de pila propio y una máquina virtual que lo ejecuta. Cada
frontera del compilador responde a una decisión de diseño justificada por las
necesidades de las fases vecinas. La extensión propia ---un sistema de macros
sintácticas con higiene y expansión recursiva, junto con \repo{match}--- aporta
metaprogramación a nivel de sintaxis sin incrementar la complejidad de las fases de
bajo nivel, gracias a su resolución como transformaciones AST$\to$AST. El resultado
es un sistema observable por capas que supera la totalidad de la batería de pruebas
de entrega final.

\end{document}
